% F19：本书原创矢量示意；依据所在节的定义、证明或明确模型。
\begin{figure}[htbp]
\centering
\begin{tikzpicture}[x=1cm,y=1cm,>=stealth,every node/.style={font=\small},line width=.65pt]

\draw[->,gray](0,0)--(5.7,0)node[right]{\(x\)};
\draw[AnalysisBlue,thick](.4,0)--(.4,1.2)--(1.7,1.2)--(1.7,2.4)--(3,2.4)--(3,1.2)--(4.3,1.2)--(4.3,0);
\draw[dashed,orange!80!black](0,.6)--(5,.6);\draw[dashed,orange!80!black](0,1.8)--(5,1.8);
\node at(2.6,3.25){阶梯函数：总变差 \(4\)};
\draw[AnalysisBlue,very thick](7.4,.6)--(11.3,.6);
\draw[AnalysisBlue,very thick](8.7,1.8)--(10,1.8);
\foreach\x/\y in{7.4/.6,11.3/.6,8.7/1.8,10/1.8}{\fill[orange!80!black](\x,\y)circle(2pt);}
\node at(9.4,3.25){两类超水平集均有周长 \(2\)};
\node at(6,-.65){\(\int_0^2P(\{u>t\})\,\mathrm{d}t=2+2=4\)};

\end{tikzpicture}
\caption[BV 的层集与周长积分]{BV 的层集与周长积分。一维紧支集阶梯模型在每个非临界层值都有两个边界点。变差把跳跃高度计入，而每层周长只计边界。}
\label{b1:fig:visual-f19}
\end{figure}
